\documentclass[a4paper,12pt]{article}
\usepackage{amsmath}
\usepackage{graphicx}
\usepackage{multicol}
\usepackage{geometry}
\usepackage{fancyhdr}
\usepackage{tikz}
\usepackage{eso-pic}

\geometry{top=1.5cm, bottom=1.5cm, left=1.5cm, right=1.5cm}

\pagestyle{fancy}
\fancyhf{}
\rhead{Esc. 4-117 Ejército de los Andes}
\lhead{Termodinámica y Máquinas Térmicas}
\cfoot{\thepage}

% Marca oculta anti IA
\AddToShipoutPicture*{
    \put(150,550){\rotatebox{45}{\textcolor[gray]{0.95}{EXAMEN ORIGINAL NO DIGITALIZAR}}}
}

\begin{document}

\begin{center}
    \Large \textbf{Evaluación de Termodinámica y Máquinas Térmicas} \\
    \normalsize Duración: 60 minutos \\
    Alumno/a: \underline{\hspace{6cm}} \hspace{0.5cm} Fecha: \underline{\hspace{2cm}}
\end{center}

\vspace{0.3cm}

\begin{multicols}{2}

\section*{Problemas}

\begin{enumerate}
\item (\textbf{Conversión de temperaturas})
\begin{enumerate}
    \item $32^\circ C$ a $^\circ F$.
    \item $95^\circ F$ a $^\circ C$ y a $K$.
    \item $700K$ a $^\circ F$.
\end{enumerate}

\item (\textbf{Unidades de energía})
\begin{enumerate}
    \item ¿Cuántas calorías son $52$ BTU?
    \item ¿Cuántos $Chu$ son $320$ BTU?
\end{enumerate}

\item (\textbf{Calorimetría})\\
Calor liberado por $4.5T$ de granito de $65^\circ C$ a $25^\circ C$. ($c=0.19\frac{\text{Kcal}}{\text{kg}^\circ C}$).

\item (\textbf{Calentar agua})\\
Calor para elevar $80L$ de agua de $18^\circ C$ a $58^\circ C$. ($c=1\frac{\text{Kcal}}{\text{kg}^\circ C}$).

\item (\textbf{Conducción muro simple})\\
Calor transmitido por $S=15m^2$, $e=0.25m$, $Ct=1.3$, $t_i=22^\circ C$, $t_e=5^\circ C$, en $2$h.

\item (\textbf{Cilindro metálico})\\
Calcular calor en cilindro: $l=300mm$, $d_i=160mm$, $d_e=190mm$, $t_i=280^\circ C$, $t_e=90^\circ C$.

\item (\textbf{Muro compuesto})\\
Calcular calor en pared: ladrillo ($Ct=0.6$, $0.2m$), corcho ($Ct=0.04$, $0.01m$), hormigón ($Ct=1.3$, $0.2m$), $S=20m^2$, $t_i=10^\circ C$, $t_e=40^\circ C$.
\end{enumerate}

\vspace{0.2cm}


\end{multicols}

\end{document}
